\chapter{Encoding Tabellen}
\label{EncodingTabellen}

\begin{figure}[h]
	\makebox[\textwidth]{\epsfxsize=\textwidth \epsffile{ISO8859-1.eps} }
	\caption{ISO Latin 1 oder auch ISO 8859-1}
	\label{isolatin1}
\end{figure}

\begin{figure}[h]
	\makebox[\textwidth]{\epsfxsize=\textwidth \epsffile{ISO8859-2.eps} }
	\caption{ISO Latin 2 oder auch ISO 8859-2}
	\label{isolatin2}
\end{figure}

\begin{figure}[h]
	\makebox[\textwidth]{\epsfxsize=\textwidth \epsffile{ISO8859-3.eps} }
	\caption{ISO Latin 3 oder auch ISO 8859-3}
	\label{isolatin3}
\end{figure}

\begin{figure}[h]
	\makebox[\textwidth]{\epsfxsize=\textwidth \epsffile{ISO8859-4.eps} }
	\caption{ISO Latin 4 oder auch ISO 8859-4}
	\label{isolatin4}
\end{figure}

\begin{figure}[h]
	\makebox[\textwidth]{\epsfxsize=\textwidth \epsffile{ISO8859-5.eps} }
	\caption{ISO Latin 5 oder auch ISO 8859-5}
	\label{isolatin5}
\end{figure}

\begin{figure}[h]
	\makebox[\textwidth]{\epsfxsize=\textwidth \epsffile{KOI8-R.eps} }
	\caption{KOI8-R}
	\label{koi8-r}
\end{figure}

